\documentclass[12pt]{article}
\usepackage[utf8]{inputenc}
\usepackage[T1]{fontenc}
\usepackage{hyperref}
\title{Saxwe Gbe Language Resources for K-12 Curriculum Development}
\author{Compiled by Claude Sonnet 5.5 (claude-sonnet-5-5) for Andr\'as Kornai}
\date{}
\begin{document}
\maketitle

\section{Introduction}
Saxwe Gbe (Saxw\textepsilon{}gbe, also Sahou\`e or Tsaphe; ISO 639-3 sxw) is a Gbe language of south-western Benin. Search results name three dialects (Saxwe, Daxe, Se). Written material is limited, and this report rests on a single web search.

\section{Primary Educational Resources}
No primer, textbook or school programme was found in the search results.

\section{Linguistic and Reference Materials}
\begin{itemize}
\item Saxwe language, Wikipedia overview: \url{https://en.wikipedia.org/wiki/Saxwe_language}
\item Saxwe (saxw\textepsilon{}gbe), Omniglot page: \url{https://www.omniglot.com/writing/saxwe.htm}
\item Webonary site for Saxwe, introduction page. Its contents were not checked, so whether it holds a dictionary is unconfirmed: \url{https://www.webonary.work/saxwe/overview/introduction/}
\end{itemize}
No published dictionary or grammar was confirmed.

\section{Religious and Community Resources}
\begin{itemize}
\item Joshua Project, Saxwe Gbe in Benin, which lists Audio Bible and Jesus Film resources: \url{https://joshuaproject.net/people_groups/14797}
\item Joshua Project language page: \url{https://joshuaproject.net/index.php/languages/sxw}
\item ScriptureEarth listing for the Gbe language code gbh. It may be a different Gbe language, so it should be checked before use: \url{https://www.scriptureearth.org/00eng.php?iso=gbh}
\end{itemize}

\section{Digital and Audio Resources}
Global Recordings Network lists recordings in Saxwe Gbe: \url{https://www.globalrecordings.net/en/language/sxw}

\section{Curriculum Development Context}
Search results gave no evidence of Saxwe being taught in schools. Literacy materials appear to be mainly Bible portions and periodicals. A native-language curriculum would need new materials, and the Webonary site and local Benin organisations should be investigated further.
\end{document}